% =====================================================================
%  Kodutöö 6 aruanne -- Algoritmid ja andmestruktuurid (ICD0001)
%  Kompileerimine: vt report/Makefile  (make  ->  aruanne.pdf)
%
%  Märkused:
%   \todo{...}  -- punane märge kohtadesse, kuhu tuleb tekst kirjutada.
%                 Enne esitamist: \renewcommand{\todo}[1]{} või kustuta.
% =====================================================================
\documentclass[12pt,a4paper,oneside]{report}

% ---------- Keel ja kodeering ----------
\usepackage[T1]{fontenc}
\usepackage[utf8]{inputenc}
\usepackage[estonian]{babel}
\usepackage{lmodern}
\usepackage{microtype}

% ---------- Lehekülje paigutus ----------
\usepackage[a4paper,left=3cm,right=2cm,top=2.5cm,bottom=2.5cm]{geometry}
\usepackage{setspace}
\onehalfspacing
\usepackage{parskip}
\usepackage{titlesec}
\titleformat{\chapter}{\normalfont\LARGE\bfseries}{\thechapter.}{0.6em}{}
\titlespacing*{\chapter}{0pt}{0pt}{1.2em}
% Peatükk ei alga uuelt lehelt
\usepackage{etoolbox}
\makeatletter
\patchcmd{\chapter}{\if@openright\cleardoublepage\else\clearpage\fi}{\par\bigskip}{}{\PackageWarning{aruanne}{chapter patch failed}}
\makeatother

% ---------- Matemaatika, tabelid, joonised ----------
\usepackage{amsmath,amssymb,amsthm}
\usepackage{booktabs,tabularx,longtable}
\usepackage{graphicx}
\graphicspath{{img/}}
\usepackage{float}
\usepackage{caption}
\usepackage{subcaption}
\usepackage{tikz}
\usetikzlibrary{arrows.meta,positioning,calc}
\usepackage{xcolor}

% ---------- Pseudokood ----------
\usepackage{algorithm}
\usepackage{algpseudocode}
\floatname{algorithm}{Algoritm}
\renewcommand{\listalgorithmname}{Algoritmide loetelu}
\algrenewcommand\algorithmicrequire{\textbf{Sisend:}}
\algrenewcommand\algorithmicensure{\textbf{Väljund:}}

% ---------- Programmikood (hele taust!) ----------
\usepackage{listings}
\definecolor{kommentaar}{RGB}{0,110,0}
\definecolor{votmesona}{RGB}{0,0,160}
\definecolor{sone}{RGB}{160,0,0}
\definecolor{taust}{RGB}{250,250,250}
\lstset{
  language=Java,
  basicstyle=\ttfamily\footnotesize,
  keywordstyle=\color{votmesona}\bfseries,
  commentstyle=\color{kommentaar},
  stringstyle=\color{sone},
  backgroundcolor=\color{taust},
  frame=single, rulecolor=\color{black!30},
  numbers=left, numberstyle=\tiny\color{black!50}, numbersep=6pt,
  breaklines=true, breakatwhitespace=false,
  showstringspaces=false, tabsize=3, columns=fullflexible, keepspaces=true,
  inputencoding=utf8, extendedchars=true,
  literate={õ}{{\~o}}1 {Õ}{{\~O}}1 {ä}{{\"a}}1 {Ä}{{\"A}}1
           {ö}{{\"o}}1 {Ö}{{\"O}}1 {ü}{{\"u}}1 {Ü}{{\"U}}1
           {š}{{\v s}}1 {Š}{{\v S}}1 {ž}{{\v z}}1 {Ž}{{\v Z}}1
}
\renewcommand{\lstlistingname}{Programmilõik}
% Tekstisisesed koodikastid ei poolitu lehtede vahel (lisades \lstinputlisting poolitub)
\BeforeBeginEnvironment{lstlisting}{\par\noindent\begin{minipage}{\linewidth}}
\AfterEndEnvironment{lstlisting}{\end{minipage}\par}

% ---------- Viited, lingid ----------
\usepackage[hidelinks]{hyperref}
\usepackage{url}

% ---------- Abikäsud ----------
\newcommand{\todo}[1]{\textcolor{red}{\textbf{[TODO:} #1\textbf{]}}}
\newcommand{\code}[1]{\texttt{#1}}
\theoremstyle{definition}
\newtheorem{definitsioon}{Definitsioon}[chapter]
\newtheorem{naide}{Näide}[chapter]
\theoremstyle{plain}
\newtheorem{vaide}{Väide}[chapter]

% ---------- Metaandmed ----------
\addto\captionsestonian{\renewcommand{\bibname}{Kasutatud allikad}}
\newcommand{\autor}{Siim Haas}
\newcommand{\pealkiri}{Vähimate ümberistumistega marsruudi leidmine ühistranspordivõrgus}
\hypersetup{pdftitle={\pealkiri}, pdfauthor={\autor}}

\begin{document}

% =====================================================================
%  TIITELLEHT
% =====================================================================
\pagenumbering{Alph}  % tiitellehe unikaalne ankur (hyperref)
% Tiitelleht
\begin{titlepage}
\thispagestyle{empty}
\begin{center}
{\large\textbf{TALLINNA TEHNIKAÜLIKOOL}}\\[0.3em]
{\large Infotehnoloogia teaduskond}\\[0.3em]
{\large IT Kolledž}

\vspace*{\fill}

{\Large Aine \emph{Algoritmid ja andmestruktuurid} (ICD0001)}\\[0.5em]
{\Large individuaalne kodutöö nr 6}\\[2em]

{\LARGE\bfseries \pealkiri}\\[1.5em]

{\large Aruanne}

\vspace*{\fill}
\end{center}

\begin{flushright}
\begin{tabular}{@{}l@{}}
Autor: \autor\\
Tudengi kood: 252641IADB\\
Õppekava: IADB17/25 IT süsteemide arendus\\[1em]
Juhendaja: Jaanus Pöial\\
\end{tabular}
\end{flushright}

\vspace{3cm}

\begin{center}
Tallinn 2026
\end{center}
\end{titlepage}


% =====================================================================
%  SISUKORD
% =====================================================================
\pagenumbering{roman}
\setcounter{page}{2}
\tableofcontents
\clearpage
\pagenumbering{arabic}

% =====================================================================
%  SISSEJUHATUS
% =====================================================================
\chapter*{Sissejuhatus}
\addcontentsline{toc}{chapter}{Sissejuhatus}

Ühistranspordis eelistavad paljud sõitjad pikemat marsruuti, kui sellel
tuleb vähem kordi ümber istuda. Aruanne kirjeldab, kuidas leida
ühistranspordivõrgus marsruut, kus ümberistumisi on võimalikult vähe.

% =====================================================================
%  1. ÜLESANDE PÜSTITUS
% =====================================================================
\chapter{Ülesande püstitus}
\label{ch:pustitus}

\section{Ülesande sõnastus}
\label{sec:sonastus}
On antud suunatud graaf, mille tipud tähistavad peatusi ja kaared tähistavad ühendusi kahe peatuse vahel.
Iga kaare infoväljas on liini tähis, mis seda ühendust teenindab (kahe peatuse vahel võib olla mitu kaart eri liinidega).
Koostada meetod, mis leiab etteantud lähtepeatusest sihtpeatusesse marsruudi, kus liinivahetusi (ümberistumisi) on võimalikult vähe.
Kui selliseid marsruute on mitu, valitakse neist kõige vähemate kaartega.
Meetod tagastab marsruudi kaarte listina (List) läbimise järjekorras.
Kui sihtpeatusesse ei saa jõuda või lähte- ja sihtpeatus on sama peatus,
tagastatakse tühi list.
Kui lähte- või sihtpeatus on puudu (\code{null}), visatakse erind
\code{IllegalArgumentException}.
Algset graafi ei muudeta.

\section{Põhimõisted}
\label{sec:moisted}
Mõiste järel on toodud ingliskeelne termin ja nimi, millega sama asi programmikoodis esineb.

\begin{itemize}
  \item \textbf{Graaf} (koodis \code{Graph}) -- Tipud ja nende vahelised ühendused ehk kaared. Selles töös näeme graafi kui ühistranspordi võrgustikku, kus on peatused (tipud) ja ühendused (kaared).
  \item \textbf{Suunatud graaf} -- Graaf, mille igal kaarel on suund: kaar algab ühest tipust ja lõpeb teises. Kaar A~$\to$~B tähendab, et peatusest A saab sõita peatusse B, kuid see ei tähenda, et peatusest B saaks sõita peatusse A. Kui liin sõidab mõlemas suunas, on graafis kaks eraldi kaart: A~$\to$~B ja B~$\to$~A.
  \item \textbf{Tipp (peatus)} (koodis \code{Vertex}) -- Tipp märgib peatust. Päris elus on peatused mõlemal pool teed või on neid suuremal ristmikul rohkem (näiteks peatus ``Vabaduse väljak''). Selles töös on üks tipp = üks peatus = näiteks ``Vabaduse väljak''.
  \item \textbf{Kaar (ühendus)} (koodis \code{Arc}) -- Peatusest väljuv ühendus teise peatusse. 1 kaar = 1 liin = 1 buss, tramm või troll. Kaare infoväljas on liini tähis (nt ``1'', ``2'', ``3'' jne).
  \item \textbf{Liin (liini tähis)} (koodis kaare väli \code{line}) -- Kaares on viide liinile, mis seda ühendust teenindab. 1 kaar = 1 liin. Kahe peatuse vahel võib olla mitu eri liiniga kaart.
  \item \textbf{Lähtepeatus} (koodis parameeter \code{from}) -- Sisend, kust marsruut algab.
  \item \textbf{Sihtpeatus} (koodis parameeter \code{to}) -- Sisend, kus marsruut lõpeb.
  \item \textbf{Marsruut} (koodis \code{List<Arc>}) -- Väljund: kaarte list, mis viib lähtepeatusest sihtpeatusesse.
  \item \textbf{Ümberistumine} (ingl \emph{transfer}, koodis loendab neid meetod \code{countTransfers}) -- Marsruudi kahel järjestikusel kaarel on erinev liini tähis, st sõitja peab vahepeatuses vahetama liini. Esimesele liinile minek marsruudi alguses ei ole ümberistumine, seega ühe liiniga marsruudil on 0 ümberistumist.
  \item \textbf{Olek} (ingl \emph{state}, koodis klass \code{State}) -- Paar (peatus, liin, millega sinna peatusesse jõuti). Lähtepeatuse olekus liini veel ei ole. Sama peatus võib otsingus esineda mitmes olekus, kui sinna jõutakse eri liinidega.
  \item \textbf{Prioriteedijärjekord} (ingl \emph{priority queue}, koodis \code{PriorityQueue}) -- Järjekord, kuhu saab elemente lisada mis tahes järjekorras, kuid mis annab välja alati kõige väiksema võtmega elemendi. Selles töös on võtmeks ümberistumiste arv ja viigi korral kaarte arv.
\end{itemize}

% =====================================================================
%  2. LAHENDUSE KIRJELDUS
% =====================================================================
\chapter{Lahenduse kirjeldus}
\label{ch:lahendus}

Ülesanne on lahendatud Java programmeerimiskeeles ühe klassina
\code{GraphTask}, mille sees on abiklassid \code{Graph}, \code{Vertex} ja
\code{Arc} (graaf, tipp ja kaar) ning otsingu jaoks vajalikud abiklassid.
Lahendus ei ole pakendatud teegiks: selle kasutamiseks tuleb vajalikud
klassid ja meetodid kopeerida oma projekti. Kasutamisjuhend on
jaotises~\ref{sec:kasutamine}.

Ülesande püstituses soovitud otsing on realiseeritud klassi
\code{GraphTask.Graph} meetodis \code{findFewestTransfersRoute}.

\section{Otsingu tööpõhimõte}

Otsingu idee on lihtne: liigume lähtepeatusest võrgus edasi nii, et
järgmisena uurime alati seda kohta, kuhu oleme jõudnud kõige vähemate
ümberistumistega (viigi korral kõige vähemate kaartega). Nii on esimene
kord, kui jõuame sihtpeatuse uurimiseni, leitud parim marsruut.

Ainult peatust meelde jätta ei piisa, oluline on ka see, millise liiniga
peatusesse jõuti, sest sellest sõltub, kas järgmine sõit on ümberistumine.
Näiteks võib peatusesse B jõuda liiniga 1 kahe kaarega ja liiniga 2 kolme
kaarega. Kui B-st sõidab sihtpeatusesse edasi ainult liin 2, on parem
pikem variant, sest siis ei pea ümber istuma. Seepärast peab otsing meeles
paari \emph{(peatus, liin, millega sinna jõuti)}; koodis on see klass
\code{State}.

Otsingu sammud:
\begin{enumerate}
  \item Alguses on järjekorras ainult lähtepeatus: 0 ümberistumist ja
        0 kaart. Esimesele liinile minek ei ole ümberistumine.
  \item Võetakse järjekorrast olek, kuhu jõuti kõige vähemate
        ümberistumistega, viigi korral kõige vähemate kaartega.
  \item Kui see on sihtpeatus, on parim marsruut leitud.
  \item Muul juhul vaadatakse läbi kõik sellest peatusest väljuvad kaared.
        Kui kaare liin erineb liinist, millega peatusesse jõuti, lisandub
        üks ümberistumine; kaarte arv suureneb alati ühe võrra. Kui
        jõutakse olekusse, kus pole varem käidud, või jõutakse sinna
        paremini kui varem, jäetakse meelde uus tulemus ja kaar, mida
        mööda tuldi, ning olek lisatakse järjekorda.
  \item Kui järjekord saab tühjaks ilma sihtpeatust leidmata, ei ole
        sihtpeatusesse võimalik jõuda ja tagastatakse tühi list.
  \item Leitud marsruut pannakse kokku, liikudes meelde jäetud kaarte
        kaudu sihtpeatusest tagasi lähtepeatusesse ja pöörates saadud
        listi ümber.
\end{enumerate}

Kõik abiandmed hoitakse meetodi sees olevates tabelites
(\code{HashMap}), seega algset graafi ei muudeta.

% =====================================================================
%  3. PROGRAMMI KASUTAMISJUHEND
% =====================================================================
\chapter{Programmi kasutamisjuhend}
\label{ch:juhend}

\section{Nõuded keskkonnale}

\begin{lstlisting}[language={},numbers=none,caption={Kasutatud tööriistade versioonid}]
$ java -version
openjdk version "25.0.5" 2026-10-20 LTS
OpenJDK Runtime Environment Temurin-25.0.5+8 (build 25.0.5+8-LTS)
OpenJDK 64-Bit Server VM Temurin-25.0.5+8 (build 25.0.5+8-LTS, mixed mode, sharing)

$ javac -version
javac 25.0.5
\end{lstlisting}

\begin{lstlisting}[language={},numbers=none,caption={Projekti struktuur}]
h6/
|-- src/main/java/GraphTask.java      programm
|-- src/test/java/GraphTaskTest.java  JUnit testid
|-- src/test/java/Aout.java           testide abiklass
|-- lib/junit-4.13.2.jar              testiteegid
|-- lib/hamcrest-core-1.3.jar
`-- bin/                              kompileeritud klassid
\end{lstlisting}

\section{Kompileerimine ja käivitamine}
\label{sec:kaivitamine}

Käsud käivitatakse projekti juurkaustas (\code{h6/}). Kompileerimine ei
trüki midagi, kui vigu pole. Programmi käivitamisel lahendatakse
mitu näiteülesannet ja trükitakse iga näite kohta graaf ning leitud
marsruut. Allpool on väljundist näidatud esimene
näide ja viimane rida, ülejäänu on asendatud märgiga \code{[...]}.

\begin{lstlisting}[language={},numbers=none,caption={Programmi kompileerimine ja käivitamine}]
$ javac -encoding UTF-8 -d bin src/main/java/*.java
$ java -cp bin GraphTask

E1
E -->
D --> aD_E_3 (D->E, line 3)
C --> aC_D_3 (C->D, line 3)
B --> aB_E_2 (B->E, line 2)
A --> aA_C_3 (A->C, line 3) aA_B_1 (A->B, line 1)

Route A -> E: [aA_C_3, aC_D_3, aD_E_3], arcs: 3, transfers: 0
[...]
Large graph: 2500 stops, 8700 arcs, route of 5 arcs with 2 transfers found in 44 ms
\end{lstlisting}

Testid kompileeritakse koos kaustas \code{lib/} olevate teekidega. Kuna
üks test käivitab ka meetodi \code{main}, trükitakse testide käigus sama
väljund uuesti; see on allpool asendatud märgiga \code{[...]}.

\begin{lstlisting}[language={},numbers=none,caption={Testide kompileerimine ja käivitamine}]
$ javac -encoding UTF-8 -d bin -cp "lib/*:bin" src/test/java/*.java
$ java -cp "bin:lib/*" org.junit.runner.JUnitCore GraphTaskTest
JUnit version 4.13.2
....
[...]
.......testLargeGraph: 15600 arcs, 38 arcs in route, 1 transfers, 47 ms
..
Time: 0.343

OK (13 tests)
\end{lstlisting}

\section{Meetodi kasutamine oma programmis}
\label{sec:kasutamine}

Kõigepealt luuakse graaf, seejärel lisatakse sellele peatused
(\code{createVertex}) ja peatuste vahelised ühendused (\code{createArc}).
Ühenduse loomisel antakse kaare id (seda võib kasutada ka nimena,
otsingus seda ei kasutata), lähtepeatus, sihtpeatus ja liini tähis. Seejärel kutsutakse välja meetod \code{findFewestTransfersRoute},
millele antakse lähte- ja sihtpeatus. Meetod tagastab marsruudi kaarte
listina; abimeetod \code{countTransfers} loendab marsruudi
ümberistumised.

Allolevas näites viib peatusest A peatusse C kaks teed: liiniga 1 läbi
peatuse B või otse liiniga 2. Kumbki ei nõua ümberistumist, seega valitakse
lühem: tagastatakse list \code{[aA\_C\_2]} ja ümberistumisi on 0.

\begin{lstlisting}[caption={Näide meetodi \code{findFewestTransfersRoute} kasutamisest},label={lst:kasutus}]
GraphTask.Graph g = new GraphTask().new Graph ("Linn");
GraphTask.Vertex a = g.createVertex ("A");
GraphTask.Vertex b = g.createVertex ("B");
GraphTask.Vertex c = g.createVertex ("C");
g.createArc ("aA_B_1", a, b, "1");   // liin 1: A -> B
g.createArc ("aB_C_1", b, c, "1");   // liin 1: B -> C
g.createArc ("aA_C_2", a, c, "2");   // liin 2: A -> C

List<GraphTask.Arc> route = g.findFewestTransfersRoute (a, c);
int transfers = GraphTask.countTransfers (route);   // 0
\end{lstlisting}

Erijuhtudel (sihtpeatusesse ei saa jõuda, lähte- ja sihtpeatus on sama
või üks neist on \code{null}) käitub meetod nii, nagu on kirjeldatud
jaotises~\ref{sec:sonastus}.

% =====================================================================
%  4. TESTIMISKAVA
% =====================================================================
\chapter{Testimiskava}
\label{ch:testid}

Programmi on testitud kahel viisil. Esiteks lahendab meetod
\code{GraphTask.run()} programmi käivitamisel kuus näiteülesannet. Viie
väikese näite puhul trükitakse graaf ja leitud marsruut, nii et tulemust
saab käsitsi kontrollida; kuues näide lahendab ülesande suurel juhuslikul
graafil ja mõõdab ainult otsingu aega. Teiseks on failis
\code{GraphTaskTest.java} JUnit 4 automaattestid, mis võrdlevad leitud
marsruuti oodatavaga ja annavad vea korral teada, milline juhtum
ebaõnnestus. Üks neist testidest käivitab ka meetodi \code{main}, seega
lahendatakse testide käigus ka kõik näiteülesanded.

Testid katavad olukorrad, kus lühim marsruut nõuab ümberistumist, aga
pikem mitte, kus võrdse ümberistumiste arvu korral otsustab kaarte arv,
kus samade peatuste vahel on mitu eri liiniga kaart ning kus on oluline,
millise liiniga vahepeatusesse jõuti. Lisaks kontrollitakse erijuhte
(sihtpeatusesse ei saa jõuda, graafis pole ühtegi kaart, lähte- ja
sihtpeatus on sama, \code{null}-argumendi korral visatakse erind
\code{IllegalArgumentException}), seda, et otsing graafi ei muuda, ja
tulemuste kokkulangevust kõiki võimalikke marsruute läbi vaatava otsinguga
400 väikesel juhuslikul graafil; suurel graafil (näiteülesandes 2500 ja
testis 3000 peatust) mõõdetakse otsingu aega.

Kõik 13 JUnit testi lõpevad edukalt. Testide täielik kood on toodud
lisas~\ref{app:testid}, kusjuures iga testi oodatav tulemus on kirjas
testi kontrolllausetes (\code{assertEquals}, \code{assertTrue} jne).
Testide kompileerimise ja käivitamise käsud on toodud
jaotises~\ref{sec:kaivitamine}.

% =====================================================================
%  KOKKUVÕTE
% =====================================================================
\chapter*{Kokkuvõte}
\addcontentsline{toc}{chapter}{Kokkuvõte}
Töö tulemusena valmis meetod \code{findFewestTransfersRoute}, mis leiab
ühistranspordivõrgus vähimate ümberistumistega marsruudi. Jätkutööna tuleks
lahendus pakendada teegiks ja nimetada mõisted ümber transpordivõrgu
terminite järgi.

% =====================================================================
%  LISAD
% =====================================================================
\appendix
\renewcommand{\thechapter}{\Alph{chapter}}
\titleformat{\chapter}{\normalfont\LARGE\bfseries}{Lisa \thechapter.}{0.6em}{}

\chapter{Programmi tekst}
\label{app:kood}
\lstinputlisting[caption={GraphTask.java}]{../src/main/java/GraphTask.java}

\chapter{JUnit testid}
\label{app:testid}
\lstinputlisting[caption={GraphTaskTest.java}]{../src/test/java/GraphTaskTest.java}

\chapter{Testide abiklass}
\label{app:aout}
\lstinputlisting[caption={Aout.java}]{../src/test/java/Aout.java}

\end{document}
